\section*{Ex.\ 1.9 --- Three Papers from ``OS History and Architecture''}
\label{p3:sec:ex19}


\subsection*{Paper 1-10 --- Dijkstra, ``The Structure of the \textit{``THE''}-Multiprogramming System''}
\label{p3:sec:the}

\textbf{Summary.}
Dijkstra reports a multiprogramming system built at Eindhoven for the Electrologica
EL~X8 by six people of on average half-time availability. The goals are modest: shorter
turnaround for short jobs, economical use of peripherals, automatic control of backing
store; it is explicitly not a multiaccess system.

Two ideas carry the design. The first decouples information from storage: segments are
named independently of pages, so a segment evicted to the drum need not return to the
page it came from and the free page with minimum latency is simply selected. The second
is the level hierarchy, a society of sequential processes at undefined speed ratios
coordinated only by \textbf{P} and \textbf{V} on semaphores, stacked strictly from
processor allocation at level~0 up to user programs, each level implementing an
abstraction that makes the one below it invisible.

The claimed payoff is verification rather than speed. Because a process can only generate
tasks for processes at lower levels, ``circularity is excluded'', and harmonious
cooperation is proved in roughly three stages, ending in the absence of what Dijkstra
calls the ``Deadly Embrace''. Testing found only trivial coding errors, one per 500
instructions, each located within ten minutes.

\textbf{Critical judgment.}
The case for layering is purely structural: the paper contains no performance data at
all, so we are never told what the hierarchy cost in memory, in context switches, or in
latency. The most revealing sentence
is the concession that testing was not yet complete ``but the resulting system is
guaranteed to be flawless''. That offers the rigour of a design in place of evidence
about an artifact, and it sits badly beside his own report a paragraph earlier: the proof
covers the synchronisation skeleton, while the errors he found were coding errors,
exactly the class it cannot reach.

The ordering also shows a crack the paper itself records. Keeping the hierarchy acyclic
required a segment fetched from the drum to be pinned in core until the requesting
process has accessed it, since otherwise ``finite tasks could be forced to generate an
infinite number of tasks for the segment controller''. Acyclicity bounds which processes
may be asked to do work, but not how often, so termination had to be rescued by a rule
living outside the hierarchy. The strict total order did not survive: in later systems
the pager needs the disk driver, which needs memory, which needs the pager.

Scope does quiet work throughout. Six people, one configuration, no multiaccess and no
user-written machine code is a setting where exhaustive testing is plausible in a way it
is not elsewhere. Dijkstra anticipates the objection and answers that ``the larger the
project, the more essential the structuring!'' That is asserted rather than argued. What
survived is the thesis beneath it: choose a structure because it makes correctness
provable.

\subsection*{Paper 1-2 --- Ritchie and Thompson, ``The UNIX Time-Sharing System''}
\label{p3:sec:unix}

\textbf{Summary.}
This describes the third version of UNIX, on the PDP-11/40 and /45. The scale is part of
the argument: 144K bytes of core of which UNIX occupies 42K, hardware costing as little
as \$40{,}000, less than two man years on the main system software, and about 40
installations since February 1971. It had been rewritten in C in 1973 at about one third
greater size, accepted for the gain in ease of modification.

The design is a small set of uniform abstractions. Ordinary files are unstructured
strings of bytes, so structure belongs to the programs that use a file and not to the
system; every device appears as a special file, under one name syntax and one protection
mechanism; all links to a file have equal status, a directory entry holding only a name
and an i-node pointer; and processes are built from \texttt{fork}, \texttt{execute},
\texttt{wait} and \texttt{pipe}. The Shell is not part of the kernel but an ordinary
swappable user program, and because the Shell rather than the command interprets
\texttt{<}, \texttt{>} and \texttt{|}, no command contains any code for redirection.
Section~9 reports 72 users, 14 maximum simultaneous, and 1800 commands a day at about 98
percent uptime.

\textbf{Critical judgment.}
The entire performance argument is section~4.1: one assembly of a 7621-line program in
35.9 seconds, divided 63.5 percent assembler execution, 16.5 percent system overhead and
20 percent disk wait. The authors then decline to interpret the figures or compare them
with any other system, saying only that they are generally satisfied. A design paper is
entitled to that stance, but the consequence is that every efficiency claim in the paper
rests on assertion.

The treatment of security is thin, and the mechanisms introduced most casually became the
classic problems. Protection is a user ID plus seven bits, with no group in this version,
plus set-user-ID and a super-user exempt from all checking. Set-user-ID is presented
purely as an elegance, available to any user on his own files without administrative
intervention. The paper never asks what happens when such a program can be induced to act
on its invoker's behalf, and decades of privilege-escalation bugs live in that gap.

``Everything is a file'' is also less total than the slogan it became, and the authors say
so where it is inconvenient: pipes are conceded to be ``not a completely general mechanism
since the pipe must be set up by a common ancestor'', which is why named pipes and then
sockets had to be added later. Section~9 also describes a research laboratory rather than
a production load, chess alone taking 5.3 percent of command CPU time, so the reliability
figures cannot support conclusions about UNIX under commercial pressure.

The Perspective section is candid about circumstance rather than method: the authors were
never faced with the need ``to satisfy someone else's requirements''. That is honest, and
it is also the part of the story that cannot be copied. What the paper does establish is
its own opening claim, that a powerful interactive system ``need not be expensive either
in equipment or in human effort''.

\subsection*{Paper 1-3 --- Engler, Kaashoek and O'Toole, ``Exokernel''}
\label{p3:sec:exokernel}

\textbf{Summary.}
The target is the fixed abstraction itself: any abstraction the kernel imposes embodies a
trade-off, and whichever way it is resolved some class of application pays. The authors
cite Cao et al., who report that application-level control over file caching alone can
reduce running time by 45 percent. Their response is to separate protection from
management: the kernel securely multiplexes the raw hardware and nothing more, while
untrusted \emph{library operating systems}, linked into the application's own address
space, implement processes, virtual memory and IPC. Changing the operating system becomes
relinking.

Three mechanisms make this safe: \emph{secure bindings}, which let the kernel perform the
expensive check once and enforce it cheaply thereafter; \emph{visible revocation}, which
makes the library a participant rather than a victim; and an \emph{abort protocol} that
breaks bindings by force when a library will not cooperate. The prototype, Aegis with
ExOS above it, is compared against Ultrix~4.2 on MIPS DECstations. Exception dispatch
takes 1.5\,$\mu$s against Ultrix's 130 and five times the best reported implementation,
protected control transfer is almost seven times the best reported, the dynamic packet
filter classifies TCP/IP headers in 1.5\,$\mu$s against 35 for MPF, and an
application-level stride scheduler takes under 100 lines of code.

\textbf{Critical judgment.}
The diagnosis is excellent and the evaluation cannot carry it. Nearly every number is a
microbenchmark of a single primitive, averaged over many repetitions, which as the authors
state themselves means the measurements ``do not consider cold start misses in the cache
or TLB, and therefore represent a best case''. A system whose entire premise is that real
applications are held back by the kernel is evaluated without a real application; the one
whole-program measurement, a $150\times150$ matrix multiplication, is dominated by the
compiler.

The comparison is also structurally unfair in a way the paper concedes in one sentence and
then walks past: ExOS does ``not offer the same level of functionality as Ultrix'', and the
authors ``do not expect these additions to cause large increases in our timing
measurements''. That expectation is exactly the claim that needed evidence, because the
absent functionality, swapping and a file system still under development, is where a
general-purpose system spends its money. Read alongside the admission that the prototype
``has no real users'', the ten-to-100$\times$ speedups are a weaker result than the number
suggests, since the baseline is doing strictly more work.

The security argument is asserted more than demonstrated: no adversary model is developed
and no attack attempted. The requirement that a packet filter must not ``lie'' and accept
packets destined for another application is answered with a trusted installer, or, on ``a
system that assumes no malicious processes'', by statically checking the filter language,
a premise that assumes away the threat the mechanism was meant to address. The paper
prefers the second route, since avoiding a central authority ``increases extensibility'',
which makes the weaker premise the load-bearing one. The
deeper unresolved tension is portability, since exposing physical names to untrusted
software makes the interface hardware-specific by construction. Their treatment of VM/370
shows it: they grant that it ``exports the ideal exokernel interface'' and separate
themselves from it on cost, that virtualising the base machine ``can be expensive'' and
``often requires additional hardware support''. That support duly arrived, which weakens
the cost objection; what survives is their second one, that hiding the real machine leaves
applications managing virtual resources counterproductively. The critique of fixed
abstractions was right and has outlived the architecture built on it.


\subsection*{References}
\label{p3:sec:refs}

\begin{enumerate}
\item[{[1-10]}] E.\,W. Dijkstra. ``The Structure of the \textit{``THE''}-Multiprogramming
  System.'' \emph{Communications of the ACM}, vol.~11, no.~5, May 1968, pp.~341--346.
  (Presented at the ACM Symposium on Operating System Principles, Gatlinburg, Tennessee,
  October 1967; also circulated as EWD196.)
\item[{[1-2]}] D.\,M. Ritchie and K. Thompson. ``The UNIX Time-Sharing System.''
  \emph{Communications of the ACM}, vol.~17, no.~7, July 1974, pp.~365--375. (Revised version
  of a paper presented at the Fourth ACM Symposium on Operating Systems Principles, Yorktown
  Heights, New York, October 1973.)
\item[{[1-3]}] D.\,R. Engler, M.\,F. Kaashoek and J. O'Toole Jr. ``Exokernel: An Operating
  System Architecture for Application-Level Resource Management.'' \emph{Proceedings of the
  Fifteenth ACM Symposium on Operating Systems Principles (SOSP~'95)}, 1995, pp.~251--266.
\end{enumerate}
